\subsubsection{Functions of \texorpdfstring{\(\mathbf u\cdot\sigv\)}{u.sigma}}
The projector pair from Eq.\ \eqref{eq:projectors}, together with the functional-calculus rule in Eq.\ \eqref{eq:functional-calculus}, gives the two-root decomposition of the unit-direction element \(\mathbf u\cdot\sigv\). The corresponding eigenvalue data are
\begin{align*}
\Pi(\pm) &= \tfrac{1}{2}(1\pm \mathbf u\cdot\sigv), \\
\Pi(\pm)^{2} &= \Pi(\pm), \\
\Pi(+)\Pi(-) &= 0, \\
\Pi(+) + \Pi(-) &= 1, \\
(\mathbf u\cdot\sigv)\Pi(\pm) &= \pm \Pi(\pm),
\end{align*}
so \(\mathbf u\cdot\sigv\) has the two eigenvalues \(+1\) and \(-1\).
Specified scalar values at those two roots determine the corresponding
projector function:
\begin{align*}
f(\mathbf u\cdot\sigv)
&= f(+1)\Pi(+) + f(-1)\Pi(-)\notag \\
&= \tfrac{1}{2}\bigl(f(+1)+f(-1)\bigr)
+ \tfrac{1}{2}\bigl(f(+1)-f(-1)\bigr)\,\mathbf u\cdot\sigv.
\end{align*}
These are the basic projector formulas for a unit-direction element. The first examples below collect the square root, logarithm branch data, rational inverse, Cayley transform, and the two exponential families for later reuse:
\begin{align*}
\sqrt{\mathbf u\cdot\sigv}
&= \epsilon(+)\Pi(+) + \epsilon(-) i\,\Pi(-), \\
\log(\mathbf u\cdot\sigv)
&= 2 i \pi\bigl(n\,\Pi(+) + (\tfrac12+m)\,\Pi(-)\bigr), \\
(z+\mathbf u\cdot\sigv)^{-1}
&= \frac{1}{z+1}\Pi(+) + \frac{1}{z-1}\Pi(-)\notag \\
&= \frac{z-\mathbf u\cdot\sigv}{z^{2}-1}, \\
\frac{i-\mathbf u\cdot\sigv}{i+\mathbf u\cdot\sigv}
&= i\,\mathbf u\cdot\sigv, \\
\exp(\theta\,\mathbf u\cdot\sigv)
&= \cosh(\theta)+\sinh(\theta)\,\mathbf u\cdot\sigv, \\
\exp(i\theta\,\mathbf u\cdot\sigv)
&= \cos(\theta)+i\sin(\theta)\,\mathbf u\cdot\sigv,
\end{align*}
where \(m,n\in\mathbb Z\) index independent logarithm branches, and the
inverse formula requires \(z\neq\pm1\). The signs
\(\epsilon(+),\epsilon(-)\in\{+1,-1\}\) are likewise independent; the conventional scalar
branches \(\sqrt{+1}=1\) and \(\sqrt{-1}=i\) correspond to
\(\epsilon(+)=\epsilon(-)=+1\). These identities follow by
applying the chosen scalar value separately on the two projector ranges.
The later exponential and logarithm subsections revisit the same formulas
with their domains stated for real and complex paravectors.

\subsubsection{Functions of X}
Let
\begin{align*}
X &= t + \mathbf r\cdot\sigv, \\
t &= \eqtext{scalar}(X)\in\mathbb R, \\
\mathbf r &= \eqtext{vector}(X)\in\mathbb R^{3}, \\
\mathbf u &:= \mathbf r/\lVert\mathbf r\rVert \qquad \text{when }\mathbf r\neq 0.
\end{align*}
If \(\mathbf r\neq 0\), write \(\mathbf r=\lVert\mathbf r\rVert \mathbf u\) with \(\mathbf u^{2}=1\). Then the minimal polynomial of \(X\) is
\begin{align*}
X^{2} - 2 t X + ( t^{2} - \mathbf r^{2} ) &= 0,\notag \\
\lambda(\pm) &= t \pm \lVert\mathbf r\rVert,\notag \\
0 &= (X-\lambda(-))(X-\lambda(+))\notag \\
&= X^{2} - \lambda(-)X\notag \\
&\quad - X\lambda(+) + \lambda(-)\lambda(+).
\end{align*}
With
\begin{equation*}
a:=\lambda(+),
\qquad
b:=\lambda(-),
\end{equation*}
the corresponding projector decomposition is
\begin{align*}
\Pi(\pm) &= \tfrac12(1\pm \mathbf u\cdot\sigv), \\
X &= a\,\Pi(+) + b\,\Pi(-), \\
\Pi(\pm)X &= \lambda(\pm)\Pi(\pm).
\end{align*}
If \(f\) is defined at the two eigenvalues, then
\begin{align*}
f(X)
&= f(a)\Pi(+) + f(b)\Pi(-)\notag \\
&= \tfrac12\bigl(f(a)+f(b)\bigr)
+ \tfrac12\bigl(f(a)-f(b)\bigr)\,\mathbf u\cdot\sigv.
\end{align*}
Writing \(X':=f(X)\), introduce its scalar and signed axial coefficients by
\begin{align*}
t' &= \tfrac12\bigl(f(a)+f(b)\bigr), \\
\rho'&:=\tfrac12\bigl(f(a)-f(b)\bigr),\\
X'&=t'+\rho'\,\mathbf u\cdot\sigv.
\end{align*}
For complex-valued \(f\), both \(t'\) and \(\rho'\) may be complex. Even
when they are real, \(\rho'\) is signed; the magnitude of the vector part is
\(\lvert\rho'\rvert\), not \(\rho'\) itself.
The projector summary may therefore be recorded as
\begin{align*}
\lambda(\pm)&=t\pm\lVert\mathbf r\rVert,\notag\\
\mathbf u&=\mathbf r/\lVert\mathbf r\rVert,\notag\\
\Pi(\pm)&=\tfrac12(1\pm\mathbf u\cdot\sigv),\notag\\
\Pi(\pm)X&=\lambda(\pm)\Pi(\pm),\notag\\
X&=\lambda(+)\Pi(+)+\lambda(-)\Pi(-),\notag\\
f(X)&=f\bigl(\lambda(+)\bigr)\Pi(+)+f\bigl(\lambda(-)\bigr)\Pi(-).
\end{align*}
The standard examples then follow immediately:
\begin{align*}
\exp(X)
&= \exp(t)\Bigl(\cosh(\lVert\mathbf r\rVert)+\sinh(\lVert\mathbf r\rVert)\,\mathbf u\cdot\sigv\Bigr), \\
\sqrt{X}
&= \epsilon(+)\sqrt{a}\,\Pi(+)+\epsilon(-)\sqrt{b}\,\Pi(-),\notag\\
&\qquad \epsilon(+),\epsilon(-)\in\{+1,-1\}, \\
\log(X) &= \log(a)\,\Pi(+) + \log(b)\,\Pi(-), \\
X^{-1}
&= \frac{\Pi(+)}{a} + \frac{\Pi(-)}{b}\notag \\
&= \frac{X^{*}}{t^{2}-\mathbf r^{2}}, \\
X^{n}
&= a^{n}\Pi(+) + b^{n}\Pi(-)\notag \\
&= \tfrac12(a^{n}+b^{n}) + \tfrac12(a^{n}-b^{n})\,\mathbf u\cdot\sigv.
\end{align*}
Here the scalar square roots and logarithms must be taken on specified
branches. A square root that is itself a real paravector exists when
\(a,b\geq0\), equivalently \(t\geq\lVert\mathbf r\rVert\); otherwise the
displayed roots are complex. The displayed logarithm requires \(a,b\neq0\); it is the real
principal logarithm when \(a,b>0\). The inverse requires \(ab\neq0\).
The power formula holds for \(n\geq0\), and also for negative integers when
\(ab\neq0\).
For two paravectors diagonalized along possibly different axes,
\begin{align*}
f(X)g(X')
&= (f(a)\Pi(+) + f(b)\Pi(-))\notag \\
&\quad \times (g(a')\Pi'(+)+g(b')\Pi'(-)),
\end{align*}
so
\begin{align*}
f(X)g(X')
&= f(a)g(a')\,\Pi(+)\Pi'(+)\notag\\
&\quad+f(a)g(b')\,\Pi(+)\Pi'(-)\notag \\
&\quad+f(b)g(a')\,\Pi(-)\Pi'(+)\notag\\
&\quad+f(b)g(b')\,\Pi(-)\Pi'(-).
\end{align*}
The mixed projector products are
\begin{align*}
\Pi(+)\Pi'(+)
&= \tfrac14(1+\mathbf u\cdot\mathbf u')\notag\\
&\quad + \tfrac14\bigl(\mathbf u'+\mathbf u+i(\mathbf u\times\mathbf u')\bigr)\cdot\sigv, \\
\Pi(+)\Pi'(-)
&= \tfrac14(1-\mathbf u\cdot\mathbf u')\notag\\
&\quad + \tfrac14\bigl(-\mathbf u'+\mathbf u-i(\mathbf u\times\mathbf u')\bigr)\cdot\sigv, \\
\Pi(-)\Pi'(+)
&= \tfrac14(1-\mathbf u\cdot\mathbf u')\notag\\
&\quad + \tfrac14\bigl(\mathbf u'-\mathbf u-i(\mathbf u\times\mathbf u')\bigr)\cdot\sigv, \\
\Pi(-)\Pi'(-)
&= \tfrac14(1+\mathbf u\cdot\mathbf u')\notag\\
&\quad + \tfrac14\bigl(-\mathbf u'-\mathbf u+i(\mathbf u\times\mathbf u')\bigr)\cdot\sigv.
\end{align*}
As a useful example, the Cayley-type transform is itself a projector exponential:
\begin{align*}
\frac{1-\mathbf r\cdot\sigv}{1+\mathbf r\cdot\sigv}
&= \exp\!\bigl(-2\,\operatorname{arctanh}(\mathbf r\cdot\sigv)\bigr)\notag \\
&= \exp\!\bigl(-2\,\operatorname{arctanh}(\lVert\mathbf r\rVert)\bigr)\Pi(+)\notag \\
&\quad + \exp\!\bigl(+2\,\operatorname{arctanh}(\lVert\mathbf r\rVert)\bigr)\Pi(-).
\end{align*}
The rational expression requires \(\lVert\mathbf r\rVert\neq1\).
For a real-valued \(\operatorname{arctanh}\) and the displayed real
exponentials, one further assumes \(\lVert\mathbf r\rVert<1\); otherwise a
complex branch must be selected.

\subsubsection{Minimal Polynomial and Projector Decomposition}
The two-root data of a complex paravector are collected here because the later function formulas, projector identities, and exponential identities all depend on this algebraic factorization; see Eqs.\ \eqref{eq:paravector-product} and \eqref{eq:functional-calculus}:
\begin{align*}
Z &= S + \mathbf V \cdot \sigv \in \C\oplus\C^{3}, \\
Z^{2} - 2 S Z + ( S^{2} - \mathbf V^{2} ) &= 0.
\end{align*}
When \(\mathbf V^{2}\neq0\), choose \(z\) with \(z^2=\mathbf V^2\) and set
\begin{align*}
\lambda(\pm)&:=S\pm z,\\
\Pi_Z(\pm)&:=\frac12\left(1\pm
\frac{\mathbf V\cdot\sigv}{z}\right).
\end{align*}
Then
\begin{align*}
( Z - \lambda (-) ) ( Z - \lambda (+) ) &= 0,\\
Z&=\lambda(+)\Pi_Z(+)+\lambda(-)\Pi_Z(-),\\
f(Z)&=f\bigl(\lambda(+)\bigr)\Pi_Z(+)
 \notag\\
&\quad+f\bigl(\lambda(-)\bigr)\Pi_Z(-)
\end{align*}
for any scalar function defined on the two eigenvalues. Changing the sign of
\(z\) merely exchanges the two labels.

The case \(\mathbf V^2=0\) must be treated separately. If
\(\mathbf V=\mathbf0\), then \(Z=S\) is scalar. If
\(\mathbf V\neq\mathbf0\), let \(N:=\mathbf V\cdot\sigv\); then \(N^2=0\),
\(Z=S+N\) is not diagonalizable, and for a function analytic at
\(S\),
\begin{equation*}
f(Z)=f(S)+f'(S)N.
\end{equation*}
This nilpotent formula is the continuous replacement for the two-projector
formula when the two roots coalesce.
